\begin{figure*}[t]
\centering
\begin{tikzpicture}[
  font=\small,
  vertex/.style={circle,draw,minimum size=7mm,inner sep=0pt},
  selected/.style={vertex,fill=blue!12,draw=blue!65!black,very thick},
  support/.style={draw,rounded corners,minimum width=40mm,
                  minimum height=8mm,fill=black!2},
  flow/.style={-{Stealth[length=2.2mm]},thick,draw=black!55}]
\node[font=\small\bfseries] at (1.6,3.0) {Vertex Cover instance};
\node[vertex] (a) at (0.2,1.8) {$a$};
\node[selected] (b) at (1.6,0.9) {$b$};
\node[vertex] (c) at (3.0,1.8) {$c$};
\draw[thick] (a)--(b)--(c);
\node[align=center] at (1.6,0.05) {$k=1$, cover $C=\{b\}$};

\draw[flow] (3.6,1.5)--(4.5,1.5);
\node[font=\small\bfseries] at (7.0,3.0) {Database instance};
\node[anchor=north] at (5.5,2.5) {\begin{tabular}{c}
  $R$\\\hline
  $a$\\
  \textcolor{blue!65!black}{$\boldsymbol{b}$}\\
  $c$
\end{tabular}};
\node[anchor=north] at (7.9,2.5) {\begin{tabular}{ccc}
  \multicolumn{3}{c}{$E(x,y,A)$}\\\hline
  $a$ & $b$ & $1$\\
  $b$ & $c$ & $1$
\end{tabular}};
\node[align=center] at (7.0,0.05) {$R$ deletable; $E$ immutable};

\draw[flow] (9.4,1.5)--(10.3,1.5);
\node[font=\small\bfseries] at (12.9,3.0) {Derivation supports};
\node[support] at (12.9,2.1)
  {$T_{ab}=\{R(a),\textcolor{blue!65!black}{R(b)}\}$};
\node[support] at (12.9,1.15)
  {$T_{bc}=\{\textcolor{blue!65!black}{R(b)},R(c)\}$};
\node[align=center] at (12.9,0.05)
  {Delete $S=\{R(b)\}$; both die};

\node[draw=black!25,fill=black!3,rounded corners,minimum width=15cm,
      align=center,inner sep=6pt] at (7.55,-0.95)
  {For COUNT and SUM, $Q(D)=2$ and $\theta=2$.\\
   $C$ covers every edge $\Longleftrightarrow$ $S$ destroys every edge derivation;
   $|C|=|S|\le k$.};
\end{tikzpicture}
\caption{A budget-preserving NP-hardness reduction from Vertex Cover for COUNT and SUM,
illustrated with two edges and $k=1$. For a general graph, each vertex
becomes a deletable $R$ tuple and each oriented edge becomes one immutable
$E$ tuple. The fixed body $R(x)\Join E(x,y,A)\Join R(y)$ gives one
support per edge. Setting $\theta=|E_G|$ forces all edge supports to be
hit. The proof extends the construction to MAX, MIN, and AVG using a
permanently surviving derivation; the same budget $k$ is retained.}
\label{fig:vertex-cover}
\Description{Three panels show the reduction. In the left panel, the graph
has edges a-b and b-c, and vertex b alone covers both. The middle panel
contains deletable R tuples a, b, c and immutable E tuples (a,b,1) and
(b,c,1). The right panel shows the supports R(a),R(b) and R(b),R(c).
Deleting the single tuple R(b) destroys both derivations, preserving the
Vertex Cover budget of one.}
\end{figure*}
